\section{Proposed Methodology}

The pipeline comprises seven stages: identity resolution, label derivation,
label validation, split construction, feature extraction, model training
with ablation, and statistical analysis. Each stage writes an artefact
consumed by the next, and each asserts its preconditions before running, so
that a stage cannot silently operate on incomplete input.

\subsection{Pipeline Overview}

Fig.~\ref{fig:pipeline} shows the end-to-end dataflow. Two corpora enter
stage A, where every image receives a photograph identifier and a patient
identifier. Stage B derives severity labels on the first corpus. Stage C
validates those labels and consolidates them to one label per photograph.
Stage D constructs patient-grouped folds and verifies that no identity key
spans a fold boundary. Stages E and F train models on the derived labels
and on expert labels respectively, using the same folds and the same code.
Stage G compares the two.

For an image $I$, the forward pass through the model is

\begin{equation}
\begin{aligned}
z    &= \phi_{\text{B0}}(I), \quad z \in \mathbb{R}^{1280}, \\
p    &= \big(\rho_{\text{nec}}, \rho_{\text{sl}}, \rho_{\text{gran}}\big)
        \in \Delta^{2}, \\
e    &= \psi(p), \quad e \in \mathbb{R}^{64}, \\
\tilde{z} &= z \odot \big(1 + \gamma \tanh(W_g e + b_g)\big), \\
\eta &= \mathcal{H}(\tilde{z}),
\end{aligned}
\label{eq:forward}
\end{equation}

where $\phi_{\text{B0}}$ is a frozen EfficientNet-B0 backbone \cite{tan2019} with its
classification head removed, $\psi$ is the Tissue Proportion Encoder,
$\Delta^{2}$ is the two-dimensional probability simplex, $\odot$ denotes
elementwise multiplication, $\gamma \in \mathbb{R}$ is a single learned
scalar, and $\mathcal{H}$ is a task-specific head. Section V-E describes
each component and Section V-F formalises them.

\subsection{Severity Label Derivation}

Because no public DFU corpus carries severity annotation, labels are
derived from tissue composition. Each image is resized to $128 \times 128$
and converted from RGB to CIE LAB, which separates luminance from
chrominance and reduces the effect of illumination differences between
sources. The wound pixel set is partitioned by $k$-means with $k = 3$,
using four initialisations. Clusters are ordered by the $L^{*}$ coordinate
of their centroids so that a given proportion index refers to the same
quantity in every image, and are assigned to necrotic, slough and
granulation tissue in ascending order of luminance, following the
convention that necrotic tissue is darkest and granulation lightest.

The tissue proportion for type $t$ is

\begin{equation}
\rho_t = \frac{|C_t|}{\sum_{j=1}^{3} |C_j|},
\qquad \sum_t \rho_t = 1,
\label{eq:props}
\end{equation}

where $C_t$ is the set of pixels assigned to cluster $t$. Severity is then
assigned by threshold rules anchored to published grading descriptions:

\begin{equation}
s =
\begin{cases}
\text{severe}   & \text{if } \rho_{\text{nec}} \geq 0.20, \\[2pt]
\text{mild}     & \text{if } \rho_{\text{gran}} \geq 0.67
                   \ \wedge\ \rho_{\text{sl}} \leq 0.27, \\[2pt]
\text{moderate} & \text{otherwise.}
\end{cases}
\label{eq:rule}
\end{equation}

Two implementation details are load-bearing. Reducing the number of
$k$-means initialisations to one shifts individual proportions by up to
0.166 relative to a ten-initialisation reference, which is the same order
as the effect being measured, so four initialisations are retained.
Ordering clusters by luminance is required because $k$-means returns
clusters in arbitrary order; without it the proportion vector is
meaningless across images.

\subsection{Label Validation Protocol}

A derived label is only useful if it recovers the quantity it claims to.
Four checks are applied, each capable of failing independently.

\emph{Expert agreement} compares derived grades against clinician review on
a stratified sample, reporting Cohen's $\kappa$ against each reviewer and
between reviewers. The inter-rater value is essential context: it bounds
what agreement any automated method could plausibly achieve.

\emph{Negative control} applies the identical derivation to images
containing no ulcer. The correct necrotic fraction for intact skin is zero,
so any severe grade is a false positive that the method cannot detect on
wound images, where no reference value exists.

\emph{Stability} groups images by source photograph and asks whether a
photograph's own augmented copies receive the same grade. Augmentation does
not change the clinical picture, so disagreement indicates that the label
is determined by something other than wound appearance.

\emph{Plausibility} compares the resulting class distribution against the
composition of clinical cohorts.

\subsection{Leakage-Free Evaluation Protocol}

Public wound corpora are frequently distributed after augmentation, so a
split over image files does not separate photographs. The protocol treats
evaluation identity and split grouping as separate problems.

Photograph identity is the union of two signals. Filenames encode the
source photograph through export and augmentation suffixes, and pixel
content is summarised by a perceptual hash made invariant to rotation and
mirroring, then compared by Hamming distance rather than equality so that
re-encoding does not split a photograph from its own copies. Neither
signal alone suffices: hashing over-splits copies whose appearance
diverges, and filenames are renumbered between exports. Two images
therefore belong to the same photograph unit if they share either
signal, computed as connected components over the union of both
relations. Implementation details are given in the released notebooks.

The split group is deliberately more conservative than the evaluation unit:
patients joined by any shared content cluster are placed in the same group,
including false-positive links, because an unnecessary merge costs less
than a leak. Folds are stratified over these groups.

Verification runs on the full image table rather than an aggregated table,
across five keys: patient, photograph, unit, group and content cluster. An
aggregated table hides patients dropped by first-match merges, and an
earlier version of this check reported success while two patients genuinely
spanned folds.

Finally, predictions are made per image and averaged to one score per
photograph before any metric is computed. Without this step a photograph
with eight augmented copies contributes eight times and one with a single
copy contributes once, which weights the test set toward heavily augmented
cases.

\subsection{TGO-Net Architecture}

Fig.~\ref{fig:arch} shows the architecture. The tissue proportions computed
during labelling are unused by a standard convolutional network, which sees
only pixels. TGO-Net supplies them as a second input stream while keeping
the backbone frozen, which keeps training feasible without sustained
accelerator access and makes every reported figure a linear probe on fixed
features rather than a tuned ceiling.

\textbf{Tissue Proportion Encoder.} The three-dimensional proportion vector
is lifted to a 64-dimensional embedding by a two-layer perceptron with GELU
activation \cite{hendrycks2016} and layer normalisation \cite{ba2016}.
Layer normalisation is used rather than batch normalisation because batch
statistics over a three-dimensional input are unstable at small batch
sizes. This component holds 2{,}432 parameters.

\textbf{Cross-Modal Gate.} The embedding is converted into a per-channel
multiplier on the frozen feature. The scalar $\gamma$ is initialised at
zero, so at initialisation $\tilde{z} = z$ and the model is exactly the
frozen-CNN baseline; the tissue pathway must earn any influence through
training. This makes $\gamma$ a directly interpretable measure of how much
the network relies on tissue composition relative to appearance. The
hyperbolic tangent bounds each channel's scaling to
$(1 - |\gamma|,\ 1 + |\gamma|)$. Multiplicative fusion is used rather than
concatenation because a 64-dimensional signal concatenated to 1{,}280
dimensions can be ignored through near-zero weights; the concatenation
variant is retained as an ablation so that this choice is tested rather
than asserted. This component holds 83{,}201 parameters.

\textbf{Prediction heads.} For the three-class ordinal formulation the head
is CORAL \cite{cao2020}: a single shared weight vector with $K-1$
independent bias terms. Because the weight vector is shared and only the
biases differ, the logits differ by a constant for every input and cannot
cross, so the model is structurally incapable of asserting
$P(y > \text{moderate}) > P(y > \text{mild})$. Rank consistency is
guaranteed by the architecture rather than encouraged by the loss. For the
binary formulation the head is a single logit.

Total trainable parameters are 86{,}915 for the ordinal head and 86{,}914
for the binary head, against 5{,}288{,}548 parameters in the frozen
backbone, so 1.62\% of the model is trained.

\subsection{Mathematical Formulation}

The Tissue Proportion Encoder maps $p \in \Delta^{2}$ to $e$ by

\begin{equation}
e = \sigma_{\text{G}}\Big(
    \mathrm{LN}\big(W_2\,
    \sigma_{\text{G}}(\mathrm{LN}(W_1 p + b_1))
    + b_2\big)\Big),
\label{eq:encoder}
\end{equation}

where $W_1 \in \mathbb{R}^{32 \times 3}$,
$W_2 \in \mathbb{R}^{64 \times 32}$, $\mathrm{LN}$ denotes layer
normalisation and $\sigma_{\text{G}}$ the GELU nonlinearity.

The Cross-Modal Gate produces a bounded per-channel modulation

\begin{equation}
g = \tanh(W_g e + b_g), \qquad
\tilde{z} = z \odot (1 + \gamma g),
\label{eq:gate}
\end{equation}

with $W_g \in \mathbb{R}^{1280 \times 64}$ and $\gamma \in \mathbb{R}$
initialised at $0$. Because $g \in (-1, 1)^{1280}$, each channel of $z$ is
scaled within $(1 - |\gamma|,\ 1 + |\gamma|)$, so a large $\gamma$ is
required before the tissue stream can materially alter the representation.

The CORAL head produces $K-1$ ordered logits from a shared projection

\begin{equation}
\eta_k = w^{\top} \tilde{z} + b_k, \qquad k = 1, \dots, K-1,
\label{eq:coral}
\end{equation}

trained against binary targets $y_k = \mathbb{1}[y > k]$ with

\begin{equation}
\mathcal{L}_{\text{CORAL}} = -\sum_{k=1}^{K-1}
\Big[ y_k \log \varsigma(\eta_k)
    + (1 - y_k) \log\big(1 - \varsigma(\eta_k)\big) \Big],
\label{eq:coralloss}
\end{equation}

where $\varsigma$ is the logistic function. The predicted class is
$\hat{y} = \sum_{k} \mathbb{1}[\varsigma(\eta_k) > 0.5]$. For the binary
task the objective is weighted cross-entropy

\begin{equation}
\mathcal{L}_{\text{BCE}} = -\big[ w^{+} y \log \varsigma(\eta)
    + (1-y) \log(1 - \varsigma(\eta)) \big],
\label{eq:bce}
\end{equation}

with $w^{+} = n_{-} / n_{+}$ computed on the training split of each fold.

Where class imbalance is addressed by resampling, SMOTE \cite{chawla2002}
is applied to the concatenated vector $[\,z \mid p\,]$ of the training
split rather than to features alone. A convex combination of two points of
$\Delta^{2}$ remains in $\Delta^{2}$, so interpolated tissue vectors stay
valid proportions:

\begin{equation}
\tilde{p} = \lambda p_a + (1-\lambda) p_b \in \Delta^{2},
\qquad \lambda \sim U(0,1).
\label{eq:smote}
\end{equation}

\subsection{Explainability}

Grad-CAM \cite{selvaraju2020} is computed on the last convolutional block.
For class score $y^{c}$ and feature maps $A^{k}$,

\begin{equation}
\alpha^{c}_{k} = \frac{1}{Z}\sum_{i}\sum_{j}
\frac{\partial y^{c}}{\partial A^{k}_{ij}},
\qquad
L^{c} = \mathrm{ReLU}\Big(\sum_{k} \alpha^{c}_{k} A^{k}\Big).
\label{eq:gradcam}
\end{equation}

The fraction of activation falling inside a tissue cluster mask $M^{c}$ is

\begin{equation}
A^{c}_{\text{align}} =
\frac{\sum_{ij} L^{c}_{ij} M^{c}_{ij}}{\sum_{ij} L^{c}_{ij}}.
\label{eq:align}
\end{equation}

Because a cluster covering a larger share of the image absorbs more
activation by area alone, we report the area-normalised lift
$A^{c}_{\text{align}} / \bar{M}^{c}$, where $\bar{M}^{c}$ is the fraction
of pixels in the mask. A lift of $1.0$ indicates attention proportional to
cluster size and therefore no preferential targeting. Section VII-E states
what this quantity can and cannot support.
